\section{Evaluation}
\label{sec:eval}

Our evaluation asks six questions.

\begin{itemize}[noitemsep,topsep=3pt,leftmargin=*]
  \item[\textbf{Q1}] Against the tools practitioners use today, how much does
        computing on the container save in time and in memory?
        (\S\ref{subsec:eval-deconstruct}--\S\ref{subsec:eval-pav})
  \item[\textbf{Q2}] Does it change what is \emph{possible}, not only what is
        fast? (\S\ref{subsec:eval-scale})
  \item[\textbf{Q3}] Is the win really about avoiding materialization, or would
        expanding to fast local storage have been enough?
        (\S\ref{subsec:eval-expand})
  \item[\textbf{Q4}] Does a better compression ratio make queries faster, as the
        design predicts? (\S\ref{subsec:eval-ratio})
  \item[\textbf{Q5}] How does the engine behave against the succinct-index
        alternative, and against its own space-time dial?
        (\S\ref{subsec:eval-gbz}--\S\ref{subsec:eval-cache})
  \item[\textbf{Q6}] Where does the approach fail? (\S\ref{subsec:eval-limits})
\end{itemize}

\subsection{Setup}
\label{subsec:eval-setup}

\paragraph{Platform.}
All runs use a single node with an AMD Ryzen Threadripper PRO 9955WX ($16$
physical cores, $32$ hardware threads), $503$\,GiB of DDR5 memory, and a Samsung
990 PRO NVMe SSD. Each command runs alone on the node, at $1$ and $16$ threads,
wrapped in \cmd{/usr/bin/time -v} to record wall-clock time and peak resident set
size. \sys reads the compressed container directly; every baseline reads the
uncompressed GFA, and \cmd{odgi} additionally builds an \cmd{.og} index first,
which we time separately and report rather than hide.

\paragraph{Cache state.}
Because the central claim of this paper concerns bytes moved, cache state is part
of the measurement rather than a detail. Unless stated otherwise, results are
\emph{warm}: the page cache is primed for every tool, which is the configuration
most favorable to the baselines, since they read far more data than \sys does.
Section~\ref{subsec:eval-expand} repeats the key comparisons \emph{cold}, dropping
caches before each run, which is the configuration a practitioner actually meets
on a first run over a freshly downloaded release.

\paragraph{Datasets.}
Table~\ref{tab:datasets} lists the graphs. Three are chromosome-scale and serve
as the regime where every baseline runs and detailed comparison is possible. Four
are whole-genome: the three successive HPRC releases, and HGSVC3, a
Minigraph--Cactus pangenome from a \emph{different} consortium. HGSVC3 does double
duty. It tests that the engine generalizes past a single construction pipeline,
and, being smaller than the largest HPRC graphs, it is the largest whole-genome
graph on which the graph-resident baselines still run at all, so it is where a
direct whole-genome comparison is possible.

\begin{table}[t]
\centering
\small
\setlength{\tabcolsep}{4pt}
\caption{Evaluation graphs. \sys reads the container; every baseline reads the
GFA.}
\label{tab:datasets}
\begin{tabular}{@{}lrrrl@{}}
\toprule
Graph & GFA & Container & Ratio & Source \\
\midrule
\emph{E. coli}  & $118$\,MB & $6.4$\,MB & $18.4\times$ & MGC, 13 strains \\
chr6            & $4.06$\,GB & $115$\,MB & $35.4\times$ & HPRC, smoothed \\
chr1            & $7.58$\,GB & $214$\,MB & $35.4\times$ & HPRC, smoothed \\
hprc-v1.1       & $48$\,GB  & $2.2$\,GB & $22.4\times$ & HPRC \\
HGSVC3          & $80.1$\,GB & $2.85$\,GB & $28.1\times$ & HGSVC \\
hprc-v2.0       & $358$\,GB & $4.3$\,GB & $83.8\times$ & HPRC \\
hprc-v2.1       & $369$\,GB & $4.5$\,GB & $82.8\times$ & HPRC \\
\bottomrule
\end{tabular}
\end{table}

The chr1 and chr6 graphs carry a linear CHM13 reference and are used for
\cmd{deconstruct} and \cmd{pav}. The \emph{E. coli} graph is circular and stores
only W-line walks, so it is not a meaningful target for reference-anchored
deconstruction nor for \cmd{odgi pav}, which does not load W lines; we retain it
as a small-graph \cmd{growth} data point.

\paragraph{Baselines.}
We compare against the reference implementation for each analysis:
\cmd{vg deconstruct}~\cite{garrison2018variation} for variant projection,
Panacus~\cite{parmigiani2024panacus} for growth curves, and
\cmd{odgi pav}~\cite{guarracino2022odgi} for presence/absence. These are the tools
the community uses, and each is the fastest available implementation of its
analysis known to us.

\paragraph{Output equivalence.}
A speedup means nothing if the answer differs, so every comparison below is
predicated on agreement, which we establish first and then stop repeating.
\cmd{growth} reproduces Panacus's curve \emph{exactly}, matching its reported node
total at every growth point on every graph. \cmd{pav} matrices are structurally
identical to \cmd{odgi}'s and reproduce its reference semantics.
\cmd{deconstruct} agrees with \cmd{vg} to within $0.13\%$ of records on the
chromosome graphs ($1{,}593{,}899$ vs.\ $1{,}593{,}956$ on chr1; $1{,}304{,}121$
vs.\ $1{,}302{,}385$ on chr6) and $0.17\%$ on HGSVC3 ($25{,}323{,}695$ vs.\
$25{,}368{,}045$), consistent with the $\sim$$99.99\%$ position concordance of
Section~\ref{sec:deconstruct}. \sys's own output is deterministic: VCFs are
byte-identical at $1$ and $16$ threads.

\subsection{Variant Projection}
\label{subsec:eval-deconstruct}

Table~\ref{tab:deconstruct} reports GFA\,$\rightarrow$\,VCF projection against
\cmd{vg deconstruct}. On the chromosome graphs \sys is $36$--$71\times$ faster
using $3.3$--$3.4\times$ less memory. The gap widens with scale, which is the
pattern to watch across all three analyses: \cmd{vg} must hold the whole graph, so
its cost tracks the GFA, while \sys's tracks the container.

\begin{table}[t]
\centering
\small
\setlength{\tabcolsep}{3.6pt}
\caption{\cmd{deconstruct} vs.\ \cmd{vg deconstruct}. ``---'' marks a run that
does not complete within node memory, or that we did not run because it was
impractical ($>$$10$\,h extrapolated).}
\label{tab:deconstruct}
\footnotesize
\begin{tabular}{@{}llrrrr@{}}
\toprule
 & & \multicolumn{2}{c}{1 thread} & \multicolumn{2}{c}{16 threads} \\
\cmidrule(lr){3-4}\cmidrule(lr){5-6}
Graph (GFA) & Tool & Time & RSS & Time & RSS \\
\midrule
\multirow{2}{*}{chr1 ($7.6$\,GB)} & \cmd{vg} & $1245$\,s & $30.1$\,GB & $805$\,s & $30.3$\,GB \\
 & \sys & $34.3$\,s & $8.8$\,GB & $11.3$\,s & $8.3$\,GB \\
\midrule
\multirow{2}{*}{chr6 ($4.1$\,GB)} & \cmd{vg} & $898$\,s & $18.4$\,GB & $349$\,s & $18.4$\,GB \\
 & \sys & $21.9$\,s & $5.7$\,GB & $7.0$\,s & $5.2$\,GB \\
\midrule
\multirow{2}{*}{HGSVC3 ($80$\,GB)} & \cmd{vg} & --- & --- & $132.9$\,m & $396.9$\,GB \\
 & \sys & $127.6$\,m & $51.7$\,GB & $8.7$\,m & $51.9$\,GB \\
\midrule
hprc-v1.1 ($48$\,GB)  & \sys & $79.9$\,m & $36.8$\,GB & $5.1$\,m & $37.4$\,GB \\
hprc-v2.0 ($358$\,GB) & \sys & --- & --- & $39.8$\,m & $191$\,GB \\
hprc-v2.1 ($369$\,GB) & \sys & --- & --- & $41.4$\,m & $195$\,GB \\
\bottomrule
\end{tabular}
\end{table}

HGSVC3 makes the residency argument concrete. \cmd{vg} completes only by consuming
$396.9$\,GB, four-fifths of the node's memory, to project an $80$\,GB graph; \sys
does the same projection $15.2\times$ faster in $51.9$\,GB. On the two largest HPRC
graphs \cmd{vg} does not complete at all, while \sys emits genome-wide VCFs of
$23$--$36$\,M records in about $40$ minutes.

This is also the analysis with the weakest memory result, and it is worth saying
why rather than burying it: deconstruction must hold per-site allele sequences and
the segment DNA payload, so peak RSS here is comparable to the uncompressed GFA
rather than far below it. The lighter analyses that follow do stay well under.

\subsection{Pangenome Growth}
\label{subsec:eval-growth}

Growth curves are the case where the disparity between the size of the question
and the cost of answering it is starkest, since the output is a few thousand
integers. Table~\ref{tab:growth} compares against Panacus.

\begin{table}[t]
\centering
\small
\setlength{\tabcolsep}{3.6pt}
\caption{\cmd{growth} vs.\ Panacus. Both tools report identical curves.
Single-threaded Panacus on the whole-genome graphs was not run.}
\label{tab:growth}
\begin{tabular}{@{}llrrrr@{}}
\toprule
 & & \multicolumn{2}{c}{1 thread} & \multicolumn{2}{c}{16 threads} \\
\cmidrule(lr){3-4}\cmidrule(lr){5-6}
Graph (GFA) & Tool & Time & RSS & Time & RSS \\
\midrule
\multirow{2}{*}{\emph{E. coli}} & Panacus & $0.98$\,s & $0.13$\,GB & $0.89$\,s & $0.13$\,GB \\
 & \sys & $0.07$\,s & $0.02$\,GB & $0.01$\,s & $0.03$\,GB \\
\midrule
\multirow{2}{*}{chr6} & Panacus & $57.3$\,s & $3.75$\,GB & $9.0$\,s & $3.77$\,GB \\
 & \sys & $3.4$\,s & $0.24$\,GB & $0.36$\,s & $0.31$\,GB \\
\midrule
\multirow{2}{*}{chr1} & Panacus & $109.9$\,s & $6.13$\,GB & $18.3$\,s & $6.16$\,GB \\
 & \sys & $7.3$\,s & $0.46$\,GB & $0.74$\,s & $0.66$\,GB \\
\midrule
\multirow{2}{*}{hprc-v1.1} & Panacus & --- & --- & $23.8$\,min & $45.8$\,GB \\
 & \sys & $69.2$\,s & $4.9$\,GB & $6.2$\,s & $6.3$\,GB \\
\midrule
\multirow{2}{*}{HGSVC3} & Panacus & --- & --- & $40.6$\,min & $71.1$\,GB \\
 & \sys & $102.9$\,s & $6.5$\,GB & $9.6$\,s & $8.1$\,GB \\
\midrule
\multirow{2}{*}{hprc-v2.0} & Panacus & --- & --- & $245$\,min & $327$\,GB \\
 & \sys & --- & --- & $39.8$\,s & $12.9$\,GB \\
\midrule
\multirow{2}{*}{hprc-v2.1} & Panacus & --- & --- & $251$\,min & $336$\,GB \\
 & \sys & --- & --- & $41.2$\,s & $13.7$\,GB \\
\bottomrule
\end{tabular}
\end{table}

On the chromosome graphs \sys is $15$--$25\times$ faster with $9$--$16\times$ less
memory, and its peak RSS stays in the hundreds of megabytes, more than $10\times$
below the uncompressed GFA. This is the sub-GFA residency thesis in its cleanest
form: the traversal store stays grammar-compressed, and all that is added is a
coverage vector.

At whole-genome scale the same mechanism produces a much larger gap, because
Panacus's cost tracks the GFA while \sys's tracks the container. Panacus needs
$245$ and $251$ minutes and $327$--$336$\,GB on the two largest HPRC graphs; \sys
computes the identical curve in $39.8$ and $41.2$ seconds using $12.9$ and
$13.7$\,GB, a $369\times$ speedup at $25\times$ less memory, with peak RSS about
$27\times$ below the GFA. The trend across the table is the point: the advantage
is not a constant factor but grows with the compression ratio of the input, which
is the prediction of \S\ref{subsec:eval-ratio}.

\subsection{Presence/Absence Variation}
\label{subsec:eval-pav}

Table~\ref{tab:pav} compares \cmd{pav} against \cmd{odgi pav}, computing the
whole-path PAV matrix ($2{,}262$ and $1{,}410$ reference windows $\times$ $46$
sample groups). At $16$ threads \sys is $299\times$ (chr1) and $613\times$ (chr6)
faster at $3.1$--$3.3\times$ less memory, and it reads the container directly,
avoiding the \cmd{odgi build} step that \cmd{odgi pav} requires ($61$\,s and
$31$\,s respectively, not included in the times shown). Single-threaded
\cmd{odgi pav} did not finish within a day; \sys's single-threaded run finishes in
under two minutes.

\begin{table}[t]
\centering
\small
\setlength{\tabcolsep}{3.6pt}
\caption{\cmd{pav} vs.\ \cmd{odgi pav}. \cmd{odgi} runs only on the chromosome
graphs, as it does not load W-line walks, and additionally needs an
\cmd{odgi build} (chr1 $61$\,s, chr6 $31$\,s at 16 threads) not counted here.}
\label{tab:pav}
\begin{tabular}{@{}llrrrr@{}}
\toprule
 & & \multicolumn{2}{c}{1 thread} & \multicolumn{2}{c}{16 threads} \\
\cmidrule(lr){3-4}\cmidrule(lr){5-6}
Graph (GFA) & Tool & Time & RSS & Time & RSS \\
\midrule
\multirow{2}{*}{chr1} & \cmd{odgi} & --- & --- & $3499$\,s & $31.9$\,GB \\
 & \sys & $106$\,s & $8.7$\,GB & $11.7$\,s & $9.5$\,GB \\
\midrule
\multirow{2}{*}{chr6} & \cmd{odgi} & --- & --- & $4759$\,s & $19.4$\,GB \\
 & \sys & $72$\,s & $5.8$\,GB & $7.8$\,s & $6.4$\,GB \\
\midrule
hprc-v1.1 & \sys & $19.3$\,min & $77$\,GB & $1.8$\,min & $77$\,GB \\
HGSVC3 & \sys & $42.4$\,min & $126.2$\,GB & $3.6$\,min & $126.2$\,GB \\
hprc-v2.0 & \sys & \multicolumn{4}{c}{terminated near $470$\,GB (exceeds node)} \\
hprc-v2.1 & \sys & \multicolumn{4}{c}{terminated near $470$\,GB (exceeds node)} \\
\bottomrule
\end{tabular}
\end{table}

The last two rows are a negative result and we discuss them in
\S\ref{subsec:eval-limits} rather than leave them to a footnote.

\subsection{Scale and Scaling}
\label{subsec:eval-scale}

\paragraph{What becomes possible.}
Across the three analyses, the whole-genome rows answer Q2. On the $358$ and
$369$\,GB graphs, \cmd{vg} and \cmd{odgi} cannot load the input on a $503$\,GiB
node at all, and Panacus completes only after four hours and a third of a terabyte
of memory. From a $4.5$\,GB container \sys answers the same questions in tens of
seconds to tens of minutes. The qualitative claim is not that these analyses got
faster; it is that two of the three moved from infeasible to routine, and that the
hardware they now need is a workstation rather than a large-memory node.

\paragraph{Thread scaling.}
Because the decode path never synchronizes (\S\ref{subsec:eng-parallel}), scaling
is limited by the fold rather than by the decoder. From $1$ to $16$ threads on
hprc-v1.1, \cmd{deconstruct} scales $15.6\times$ ($79.9 \to 5.1$\,min),
\cmd{pav} $10.7\times$ ($19.3 \to 1.8$\,min), and \cmd{growth} $11.2\times$
($69.2 \to 6.2$\,s). \cmd{deconstruct} scales nearly linearly because its parallel
region is a pure per-haplotype fold into private buffers; the other two lose some
of it to their shared-accumulator merges, which is the expected cost of the
strategies of \S\ref{subsec:eng-parallel}.
\needsexp{full 1/2/4/8/16/32-thread curves for all three analyses on hprc-v1.1
and chr1, to replace these endpoints with a scaling figure}

\subsection{Would Expanding to Fast Storage Have Been Enough?}
\label{subsec:eval-expand}

This is the alternative a storage practitioner reaches for first
(\S\ref{subsec:bg-alternatives}), and the comparisons above do not settle it,
because they run warm and give the baselines their input already in cache. Here we
test it directly. For each graph we measure three configurations end to end from a
cold page cache: (\emph{i}) \sys reading the container; (\emph{ii}) the baseline
reading a GFA already resident on local NVMe; and (\emph{iii}) the honest
archival workflow, in which only the compressed form is stored and the baseline's
input must first be expanded to NVMe, so the expansion time and the write
amplification count against it.

\needsexp{cold-cache end-to-end table: for chr1, hprc-v1.1, HGSVC3 report, per
configuration, wall-clock time, bytes read from the block device (from
\cmd{/proc/<pid>/io} or \cmd{iostat}), bytes written, and peak RSS. This is the
single most important missing experiment: it converts the paper's argument from a
speedup claim into a storage claim, and reviewers will ask for it explicitly.}

The expected shape of the result follows from the decomposition in
\S\ref{subsec:bg-tax}, and is worth stating as a prediction the experiment can
falsify. Expanding to NVMe removes the I/O component and nothing else; on HGSVC3
it must additionally write $80$\,GB and hold it. The parse-and-build and residency
components are untouched, and they are the terms that make \cmd{vg} take
$132.9$ minutes and $396.9$\,GB. If the measurement shows configuration
(\emph{iii}) landing close to (\emph{ii}), the storage tier was never the binding
constraint, and the paper's argument holds; if it shows the two far apart, then a
substantial share of the reported advantage is I/O that faster storage could have
bought, and we should say so.

\subsection{Does Ratio Buy Speed?}
\label{subsec:eval-ratio}

The design predicts an inversion of the usual trade-off: because the compressed
columns are what the inner loop walks, a graph that compresses better should be
queried faster per unit of biological content, rather than more slowly. The tables
above are consistent with this, since the largest speedups occur on hprc-v2.0/v2.1
where the ratio is $83\times$, and the smallest on \emph{E. coli} where it is
$18\times$. But those graphs differ in many ways besides ratio, so the comparison
is suggestive rather than controlled.

\needsexp{controlled ratio sweep: hold the graph fixed and vary the number of
grammar rounds (and hence the container's ratio) over its useful range, then plot
query wall-clock time and bytes touched for \cmd{growth} and \cmd{deconstruct}
against the resulting ratio. This isolates ratio from graph identity and directly
tests the paper's central mechanism. Expect a decreasing curve that flattens once
decode work, rather than bytes touched, dominates; the location of that knee is
itself a result worth reporting.}

\subsection{Against a Succinct Index}
\label{subsec:eval-gbz}

GBZ~\cite{siren2022gbz} is the closest existing design (\S\ref{subsec:bg-alternatives})
and deserves a direct measurement rather than an argument.

\needsexp{GBZ comparison. (a) Container size vs.\ GBZ size on every graph, which
is partly available from the compressor's prior evaluation but should be reported
here from our own runs. (b) A common analysis run over both: the natural candidate
is \cmd{vg deconstruct} operating on a GBZ rather than on GFA text, which isolates
``compressed-index-resident'' from ``text-resident'' and is the fairest possible
version of the baseline. Report GBZ construction time separately, as we do for
\cmd{odgi build}. Expect GBZ to close much of \cmd{vg}'s memory gap and little of
its time gap, which would be the cleanest evidence for the navigational-vs-fold
distinction the paper draws.}

\subsection{The Space-Time Dial}
\label{subsec:eval-cache}

The rule-leaf cache budget (\S\ref{subsec:eng-memory}) is the engine's one
explicit knob, interpolating between fully lazy expansion and full
materialization.

\needsexp{cache sensitivity sweep: for \cmd{pav} (multi-pass, the analysis that
benefits) and \cmd{growth} (single-pass, which should not), sweep the budget over
roughly $0$, $1$, $4$, $16$, $64$\,GB on chr1 and hprc-v1.1 and plot wall-clock
time and peak RSS. The expected result is a knee at modest budget for \cmd{pav}
and a flat line for \cmd{growth}; the flat line is as important as the knee,
because it shows the default of zero is right for streaming folds.}

\subsection{Where It Fails}
\label{subsec:eval-limits}

The whole-genome \cmd{pav} rows in Table~\ref{tab:pav} are the clearest limit of
the design, and they fail for a reason worth stating precisely, because it is not
the reason a reader might assume.

The traversal store is not what runs out of memory. \cmd{pav} builds a node-to-group
CSR index and per-window accumulators, and those are ordinary uncompressed
structures whose size grows with the node count and the window count, not with the
query. On HGSVC3 that working set is already $126$\,GB against an $80$\,GB input,
so this is the one analysis that does not meet the sub-GFA residency goal even
where it completes. On the two largest HPRC graphs ($94{,}554$ and $60{,}807$
windows over $148$--$166$\,M nodes) it passes $470$\,GB and is terminated.

The general statement is that computing over compression bounds the size of the
\emph{data} a query holds, and says nothing about the size of the \emph{state} the
query chooses to build. When the accumulator is $O(\text{nodes})$ or
$O(\text{windows})$, as for four of our six analyses, the bound is what matters
and residency stays far below the input. When the accumulator is a join index that
scales with the graph, the compressed store is no longer the dominant term and the
engine inherits the same memory wall as everyone else. Compressing the join index
is the obvious next step and we return to it in \S\ref{sec:discussion}.
